\paragraph{Energy-Based Models}
Energy-Based Models (EBMs) represent a family of generative models that learn an unnormalised probability distribution over data by associating each sample $\mathbf{x}$ with a scalar energy value $E_\theta(\mathbf{x})$, parameterised by a neural network. 
Lower energy corresponds to higher probability, and the model distribution is defined as:
\begin{equation}
    p_\theta(\mathbf{x}) = \frac{\exp(-E_\theta(\mathbf{x}))}{Z_\theta}, \qquad Z_\theta = \int \exp(-E_\theta(\mathbf{x}))\, d\mathbf{x},
\end{equation}
where $Z_\theta$ is the intractable partition function required for normalisation.

Following \citet{NEURIPS2019_378a063b}, EBMs are trained by adjusting $E_\theta(\mathbf{x})$ using contrastive divergence \citep{product_of_experts} such that real data samples are assigned lower energies than generated or negative samples. 
The objective is derived from the maximum likelihood principle, whose gradient can be expressed as \citep{NEURIPS2020_49856ed4}:
\begin{equation}
    \nabla_\theta \mathcal{L} = 
    \mathbb{E}_{\mathbf{x}^+ \sim p_{\text{data}}}[\nabla_\theta E_\theta(\mathbf{x}^+)]
    - \mathbb{E}_{\mathbf{x}^- \sim p_\theta}[\nabla_\theta E_\theta(\mathbf{x}^-)],
\end{equation}
where $\mathbf{x}^+$ are real data points and $\mathbf{x}^-$ are samples drawn from the current model distribution $p_\theta$.

Sampling from $p_\theta$ is typically performed using Markov Chain Monte Carlo (MCMC) methods such as Langevin dynamics \citep{bayesian_learning_sgld}, which iteratively refine noisy samples according to:
\begin{equation}
    \mathbf{x}^{k} = \mathbf{x}^{k-1} - \frac{\lambda}{2}\nabla_{\mathbf{x}} E_\theta(\mathbf{x}^{k-1}) + \omega^{k}, \qquad \omega^{k} \sim \mathcal{N}(0, \lambda),
\end{equation}
where $k$ denotes the iteration step and $\lambda$ is the step size. 
This process encourages samples to move toward regions of low energy, effectively concentrating probability mass around modes of the data distribution.

EBMs are therefore an important bridge between traditional likelihood-based models and score-based diffusion frameworks. 
Their formulation naturally connects to stochastic differential equations (SDEs) and score matching, where $\nabla_\mathbf{x} \log p_\theta(\mathbf{x}) = -\nabla_\mathbf{x} E_\theta(\mathbf{x})$, enabling a seamless interpretation of diffusion models as energy-based systems. 
In this unified view, the energy function defines the geometry of the data manifold, guiding sample trajectories toward high-probability regions in the same way that score fields govern the dynamics of diffusion processes.
